\documentclass[11pt]{article}
\usepackage{amsmath,amssymb,amsthm}
\newtheorem{theorem}{Theorem}
\begin{document}
\title{An asymptotic obstruction to cyclic saturation of second-order AC bounds}
\author{}\date{}\maketitle

The following excludes all upper-dimension subspaces invariant under a cyclic
rotation of order $k$, about any axis. It does not exclude arbitrary
anticoherent subspaces attaining the general upper bound.

\begin{theorem}
Let $N\ge420$ and $k=\lfloor N/5\rfloor+1$. There is no second-order
anticoherent $k$-dimensional subspace invariant under the collective rotation
$\exp(-2\pi i J_{\mathbf n}/k)$ about any axis $\mathbf n$.
For the subsequence $N=5k-1$, the same conclusion holds already for $k\ge62$.
\end{theorem}

\begin{proof}
A collective rotation preserves anticoherence, so we may take the symmetry
axis to be $z$. Write $j=N/2$ and $c=j(j+1)/3$.

We first prove that a second-order AC subspace supported on at most five
$J_z$ eigenvalues spaced by $k$ has dimension at most one if $c>2k^2$.
Pad missing endpoints by zero coefficients, and write the five positions as
$d-2k,d-k,d,d+k,d+2k$. If the subspace had dimension at least two, the kernel
of the top-coordinate functional would contain a nonzero vector. After
normalization this vector has mean zero and variance $c$, supported in
$[d-2k,d+k]$. For any real random variable in $[L,U]$,
$(m-L)(U-m)\ge0$ implies
$\mathbb E[m^2]\le(L+U)\mathbb E[m]-LU$. Consequently
\[
 c\le 2k^2+kd-d^2.
\]
Applying the same argument to a possibly different vector in the kernel of
the bottom-coordinate functional gives
\[
 c\le2k^2-kd-d^2.
\]
Together these imply $c\le2k^2-k|d|-d^2\le2k^2$, a contradiction.

Write $N=5k-s$, $1\le s\le5$. The cyclic rotation eigenspaces are precisely
the Dicke residue sectors $r\equiv a\pmod k$, each with at most five nodes.
Moreover,
\[
 c-2k^2=\frac{k^2-10(s-1)k+s(s-2)}{12}>0
\]
for $k\ge85$, as in the stated range. For $s=1$ this expression is already
positive for every $k>1$. A cyclic-invariant code decomposes into subspaces
of these sectors; each is itself AC and thus has dimension at most one.
Since the total dimension is $k$, all $k$ sectors have dimension exactly one.
We therefore obtain an orthonormal basis $u_0,\ldots,u_{k-1}$ satisfying
\[
 \operatorname{supp}(u_a)\subseteq
 \{r:0\le r\le N,\ r\equiv a\pmod k\}.
\]
For half-integer spin, the overall phase of the rotation does not alter
these distinct residue eigenspaces.

Let $P$ be the code projector. AC order two implies
\[
 PJ_\alpha P=0,\qquad PJ_\alpha J_\beta P=c\delta_{\alpha\beta}P.
\]
Consequently the four isometries $P,J_xP/\sqrt c,J_yP/\sqrt c,J_zP/\sqrt c$
have mutually orthogonal ranges. Their range projector is
\[
 E=P+\frac1c\sum_{\alpha=x,y,z}J_\alpha PJ_\alpha\le I.
\]
Let $p_r=P_{rr}$ in the Dicke basis, and put $m_r=j-r$. Its diagonal
inequalities read
\[
 \left(1+\frac{m_r^2}{c}\right)p_r+
 \frac{r(N-r+1)p_{r-1}+(r+1)(N-r)p_{r+1}}{2c}\le1,
 \tag{1}
\]
where missing endpoint terms are zero. The support hypothesis supplies,
for each residue $a$,
\[
 \sum_{r\equiv a}p_r=1,\qquad
 \sum_{r\equiv a}m_rp_r=0,\qquad
 \sum_{r\equiv a}m_r^2p_r=c. \tag{2}
\]

Set $\theta=\pi/8$ and use seven positive averaging weights
\[
 f_a=\sin((a+1)\theta),\quad 0\le a\le6,\qquad F=\sum_{a=0}^6 f_a.
\]
Set $f_{-1}=f_7=0$. In each of the five bands $r=\ell k+a$, $0\le\ell\le4$,
sum (1) with multipliers $f_a$. All these indices exist: writing
$N=5k-s$, $1\le s\le5$, it suffices that $k\ge s+6$, which is certainly true
in the stated range.

The coefficient of an internal $p_r$ in the weighted sum is
\begin{align*}
 &\left(1+\frac{m_r^2}{c}\right)f_a+
 \frac{(j(j+1)-m_r^2+m_r)f_{a+1}
       +(j(j+1)-m_r^2-m_r)f_{a-1}}{2c}\\
 &=\left[1+3\cos\theta+
       \frac{m_r^2}{c}(1-\cos\theta)\right]f_a
       +\frac{m_r}{c}\sin\theta\cos((a+1)\theta)\\
 &\ge\left[1+3\cos\theta-\frac{6}{N+2}\right]f_a.
\end{align*}
Here $\sin\theta\,|\cot((a+1)\theta)|\le1$ and $|m_r|/c\le6/(N+2)$.
Contributions from outside the seven-index window are nonnegative.
Moreover, $\cos(\pi/8)>11/12$: squaring reduces this to
$\sqrt2>49/36$, itself an exact rational squared inequality.
Thus, with
\[
 B(N)=\frac{15}{4}-\frac6{N+2},\qquad
 W_\ell=\frac1F\sum_{a=0}^6f_a p_{\ell k+a},
\]
we obtain
\[
 0\le W_\ell\le\frac1{B(N)},\qquad \sum_{\ell=0}^4W_\ell=1. \tag{3}
\]

Define ideal normalized latitudes and a quadratic dual polynomial by
\[
 z_\ell=1-\frac{2\ell}{5},\qquad Q(z)=z^2-\frac85z.
\]
The actual normalized latitude is
\[
 \frac{m_{\ell k+a}}j=z_\ell-e_{\ell a},\qquad
 e_{\ell a}=\frac{2(a+s\ell/5)}N\le\frac{12+2\ell}{N}=:\epsilon_\ell.
\]
The five values of $Q(z_\ell)+3/5$ are
\[
 0,\quad0,\quad\frac8{25},\quad\frac{24}{25},\quad\frac{48}{25}.
\]
For band zero, $Q(1-e)\le-3/5$ whenever $0\le e\le2/5$, as holds here.
For bands $\ell=1,\ldots,4$,
\[
 Q(z_\ell-e)\le Q(z_\ell)
       +\left(\frac85-2z_\ell\right)\epsilon_\ell+\epsilon_\ell^2.
\]
All four excesses above $-3/5$ are nonnegative. Applying (3), their total
weighted expectation is bounded by
\[
 U(N)=-\frac35+
 \frac{\displaystyle\frac{16}{5}+\frac{584}{5N}+\frac{1176}{N^2}}
      {\displaystyle\frac{15}{4}-\frac6{N+2}}. \tag{4}
\]
On the other hand, the exact moments (2) imply that the same expectation is
\[
 \frac{c}{j^2}=\frac{N+2}{3N}>\frac13. \tag{5}
\]
The function $U(N)$ decreases for $N\ge420$, since its positive numerator
decreases and its positive denominator increases. Exact rational arithmetic
gives
\[
 \frac13-U(420)=\frac{872}{1655325}>0.
\]
Equations (4) and (5) contradict each other.

For $s=1$, sharpen the band error to
$\epsilon_\ell=(12+2\ell/5)/N$. The numerator in (4) becomes
\[
 \frac{16}{5}+\frac{424}{5N}+\frac{3384}{5N^2}.
\]
At $N=309=5\cdot62-1$, this gives an upper bound strictly below $1/3$;
the same monotonicity argument covers every larger $k$.
\end{proof}

The proof makes no continuity or smoothness assumption on coefficients or
phases. The one-dimension-per-residue property is a consequence of cyclic
invariance and the variance lemma, not an extra hypothesis. The result
strengthens the earlier smooth-profile obstruction, while leaving open
upper-dimension subspaces without this cyclic invariance.
The accompanying \texttt{construction\_cyclic\_exclusion.py} checks every
rational coefficient, endpoint gap, and finite-window formula used above.
\end{document}
